% บทที่ 1 บทนำ
\section{ข้อมูลทั่วไปของระบบ}
Pacegasus เป็นแอปพลิเคชันผู้ช่วยออกแบบ จัดตาราง และปรับแผนการวิ่งเฉพาะบุคคล
คู่มือฉบับนี้อธิบายการเตรียมระบบ การติดตั้ง และการใช้งานแอปพลิเคชันตามลำดับ
เพื่อให้ผู้ใช้งานเข้าใจวิธีใช้งานและสามารถปฏิบัติตามได้

\section{วัตถุประสงค์ของคู่มือ}
\begin{enumerate}[label=\arabic*)]
    \item เพื่ออธิบายการเตรียมความพร้อมและการติดตั้งระบบ Pacegasus
    \item เพื่ออธิบายขั้นตอนการใช้งานฟังก์ชันหลักของแอปพลิเคชัน
    \item เพื่อใช้เป็นเอกสารประกอบการใช้งานและตรวจสอบขั้นตอนเมื่อพบปัญหา
\end{enumerate}

\section{ขอบเขตของคู่มือ}
คู่มือแบ่งเนื้อหาเป็นการติดตั้งระบบและการใช้งานแอปพลิเคชัน
ครอบคลุมการเตรียมเครื่องมือ การตั้งค่าบริการที่เกี่ยวข้อง การเข้าสู่ระบบ
และการออกแบบ จัดตาราง และปรับแผนการวิ่งเฉพาะบุคคล

\section{กลุ่มผู้ใช้งานคู่มือ}
ผู้ใช้งานระบบแบ่งเป็น 2 กลุ่ม ดังนี้
\begin{enumerate}[label=\arabic*)]
    \item ผู้ใช้ทั่วไป ใช้งานฟังก์ชันสำหรับผู้ใช้ตามคู่มือในบทที่ 2
    \item หัวหน้าคลับ ใช้งานส่วนการจัดการคลับตามสิทธิ์ที่ระบบกำหนด
\end{enumerate}
สำหรับผู้ติดตั้งหรือผู้พัฒนา สามารถศึกษาขั้นตอนเตรียมเครื่องมือและตั้งค่าระบบได้ในบทที่ 3

\section{ข้อกำหนดเบื้องต้นในการใช้งาน}
\subsection{ฮาร์ดแวร์}
เตรียมคอมพิวเตอร์สำหรับติดตั้งและตั้งค่าระบบ รวมถึงอุปกรณ์สำหรับทดสอบแอปพลิเคชัน
% TODO: ระบุรุ่นอุปกรณ์ ระบบปฏิบัติการ และข้อกำหนดขั้นต่ำตามผลทดสอบจริง

\subsection{ซอฟต์แวร์และบริการที่เกี่ยวข้อง}
เตรียมเครื่องมือสำหรับพัฒนาแอปพลิเคชัน Flutter
และบัญชีสำหรับตั้งค่าบริการ Supabase, Google Authentication และ API แผนที่ Sphere GISTDA
โดยมีรายละเอียดการติดตั้งและตั้งค่าในบทที่ 3
% TODO: ระบุเวอร์ชันเครื่องมือและบริการที่ใช้จริง
